\subsection*{Question 2.4 - Nonlinear state-space models}

We want the state equations for the two alternative state vectors
\begin{equation*}
x = \begin{bmatrix} {}^B p \\ v \\ \lambda \\ \omega \end{bmatrix}, \qquad
x_1 = \begin{bmatrix} p \\ v \\ \lambda \\ \omega \end{bmatrix}, \qquad
{}^B p = R^T(\lambda)\, p,
\end{equation*}
with the input $u = \begin{bmatrix} T \\ n \end{bmatrix}$ from Question 2.3.
(The handout writes $\dot p$ in the second position of $x_1$, which has to be a typo, since $\dot p = R v$ is not an independent state. We take $v$ there, so $x_1 \in \mathbb{R}^{12}$.)
Three blocks are common to both models, so we derive those first and then treat the position block of each model separately.

\paragraph{Step 1 - Translational dynamics.}
The first equation in (1) is $m\dot v = -m S(\omega) v + f$.
From Question 2.1 the force in the body frame is the gravity term plus the thrust term,
\begin{equation*}
f = m g\, R^T(\lambda)\, e_3 - T e_3, \qquad
R^T(\lambda) e_3 = \begin{bmatrix} -\sin\theta \\ \sin\phi\cos\theta \\ \cos\phi\cos\theta \end{bmatrix},
\end{equation*}
where $R^T(\lambda) e_3$ is simply the third row of $R(\lambda)$ written as a column.
Dividing by $m$,
\begin{equation*}
\dot v = -S(\omega) v + g\, R^T(\lambda)\, e_3 - \frac{T}{m}\, e_3 .
\end{equation*}

\paragraph{Step 2 - Rotational dynamics and Euler angle kinematics.}
The second equation in (1), solved for $\dot\omega$ with the torque $n = u_{2:4}$, gives
\begin{equation*}
\dot\omega = J^{-1}\big(n - S(\omega) J \omega\big).
\end{equation*}
The attitude kinematics are given directly by the handout,
\begin{equation*}
\dot\lambda = Q(\lambda)\,\omega, \qquad
Q(\lambda) = \begin{bmatrix} 1 & \sin\phi\tan\theta & \cos\phi\tan\theta \\ 0 & \cos\phi & -\sin\phi \\ 0 & \frac{\sin\phi}{\cos\theta} & \frac{\cos\phi}{\cos\theta} \end{bmatrix}.
\end{equation*}

\paragraph{Step 3 - Position block for $x_1$.}
For the state $x_1$ the position is the inertial position $p$, and its derivative is the linear kinematics from the handout, $\dot p = R(\lambda)\, v$ (the velocity is stored in $\{B\}$, so it has to be rotated to $\{I\}$).
Together with Steps 1 and 2,
\begin{equation*}
\dot x_1 = g_1(x_1, u) =
\begin{bmatrix}
R(\lambda)\, v \\[4pt]
-S(\omega) v + g\, R^T(\lambda)\, e_3 - \dfrac{T}{m}\, e_3 \\[8pt]
Q(\lambda)\, \omega \\[4pt]
J^{-1}\big(n - S(\omega) J \omega\big)
\end{bmatrix}.
\end{equation*}

\paragraph{Step 4 - Position block for $x$.}
Now the position is ${}^B p = R^T p$, and both factors depend on time, so the product rule gives
\begin{equation*}
{}^B\dot p = \dot R^T p + R^T \dot p .
\end{equation*}
We need $\dot R$.
The angular velocity $\omega$ of the body, written in $\{B\}$, is related to the rotation matrix by $\dot R = R\, S(\omega)$.
This is where the hint $R(a \times b) = (Ra) \times (Rb)$ is used: it says $R\,S(a)\,b = S(Ra)\,R\,b$ for every $b$, that is, $R\,S(a) = S(Ra)\,R$.
So $R\,S(\omega)$ is the same as $S(R\omega)\,R$, the familiar inertial-frame form of $\dot R$ with the angular velocity rotated into $\{I\}$, and both forms are consistent.

Using $S(\cdot)^T = -S(\cdot)$ (the second hint),
\begin{equation*}
\dot R^T = \big(R\,S(\omega)\big)^T = S(\omega)^T R^T = -S(\omega)\, R^T .
\end{equation*}
The second term is easy: $R^T \dot p = R^T R\, v = v$.
Substituting both,
\begin{equation*}
{}^B\dot p = -S(\omega)\, R^T p + v = v - S(\omega)\, {}^B p = v - \omega \times {}^B p .
\end{equation*}
So the position derivative in the body frame is the velocity minus a term $\omega \times {}^B p$ due to the frame rotating with the vehicle.
The block is attitude-independent in the sense that it no longer contains $R(\lambda)$, which is what makes the later linearization simpler.

\paragraph{Step 5 - Final result.}
Putting the blocks together,
\begin{equation*}
\dot x = g(x, u) =
\begin{bmatrix}
v - S(\omega)\, {}^B p \\[4pt]
-S(\omega) v + g\, R^T(\lambda)\, e_3 - \dfrac{T}{m}\, e_3 \\[8pt]
Q(\lambda)\, \omega \\[4pt]
J^{-1}\big(n - S(\omega) J \omega\big)
\end{bmatrix}.
\end{equation*}
The two models $g$ and $g_1$ differ only in their first block: $v - S(\omega)\,{}^B p$ for the first and $R(\lambda)\,v$ for the second.
